\documentclass[12pt]{article}
\usepackage[utf8]{inputenc}
\usepackage[T1]{fontenc}
\usepackage{hyperref}
\title{Maay Language Resources for K-12 Curriculum Development}
\author{Compiled by Claude Sonnet 5.5 (claude-sonnet-5-5) for Andr\'as Kornai}
\date{}
\begin{document}
\maketitle

\section{Introduction}
Maay (Af-Maay, Maay Maay, Mai Mai; ISO 639-3 ymm) is a Somali-related language of southern Somalia, with some speakers in neighbouring Ethiopia and Kenya. It is partly mutually comprehensible with Northern Somali, and its speakers often use Standard Somali as a lingua franca. Speaker estimates vary by source. This report rests on a single web search.

\section{Primary Educational Resources}
\begin{itemize}
\item Af Maay online course, a commercial product (14 books, a dictionary, recorded lectures and video clips). Its quality was not verified: \url{https://afmaay.gumroad.com/l/vielki}
\end{itemize}
No school textbooks or national curriculum materials in Maay were found. Wikipedia states that Maay is not generally used in education or media.

\section{Linguistic and Reference Materials}
\begin{itemize}
\item Maay Maay (Wikipedia): \url{https://en.wikipedia.org/wiki/Maay_Maay}
\item Omniglot page on Maay Maay (Af-Mai-Mai), which notes that several Latin-script orthographies have been developed since the 1980s: \url{https://omniglot.com/writing/maaymaay.htm}
\item English-Maay Dictionary by Mohamed Haji Mukhtar, a retail listing describing over 15,000 words: \url{https://www.elefant.ro/english-maay-dictionary-hardcover_c14520cd-ec7c-41d9-9255-c2e1572dc2d0}
\item Bilinguistics page on Maay Maay speech and language development: \url{https://bilinguistics.com/languages/maay-maay-speech-and-language-development/}
\end{itemize}

\section{Religious and Community Resources}
No religious or community organisation resources were found in the search results.

\section{Digital and Audio Resources}
Only the commercial course above includes audio or video. No free corpus was found.

\section{Curriculum Development Context}
The search found no evidence of Maay-medium schooling. Several competing spellings mean an orthography decision would be an early step. The Mukhtar dictionary and the Omniglot summary are the starting points found here.
\end{document}